\begin{progprob}
    % Write the problem below. Precede any paths with \parentdir. For example:
    %   \includegraphics{\parentdir/public/fig/picture.png}
    % Any files used in the solution should be placed in the `private` 
    % subdirectory. Any other files can go in the `public` subdirectory.

    Recall that a \emph{mode} of a collection is an element which occurs with the greatest
    frequency. For example, 4 is a mode of the collection $4, 5, 8, 3, 4, 2, 4, 5, 5, -2$.
    5 is also a mode, since it occurs just as frequently as 4.

    In a file named \python{mode.py}, create a function named
    \python{mode(numbers)} which takes in an array of numbers and returns a
    mode. To receive full credit, your function should have the best possible
    average case time complexity. If there are multiple modes, your function need only
    return one of them. You should \emph{not} assume that the numbers are
    integers or that they are positive.

    \emph{Note}: it's up to you to determine the best possible time complexity and convince
    yourself that your code is efficient!
    Think back to when we talked about theoretical lower bounds. Can you come up with one
    for this problem? This is a situation you'll find yourself in often when writing code:
    there's usually no resource that tells you what the best time complexity is for a given
    problem, so you'll have to convince yourself that your solution is optimal.

    \begin{soln}
        \inputminted{python}{\thisdir/include-solution/mode.py}
    \end{soln}


\end{progprob}
